% !TeX root = ../../ONMIPA Alin.tex
% Rujukan: Zhang (2011), bagian 1.2, 2.1, 2.3; Zhang (2009),
% soal-soal tentang ruang kolom, kernel, rank, nulitas, dan matriks khusus.
\chapter{Rank dan Nulitas}
\label{ch:rank}

Rank mengukur dimensi ruang hasil suatu matriks, sedangkan nulitas
mengukur dimensi ruang masukan yang dipetakan ke nol. Keduanya
menghubungkan kebebasan kolom, banyaknya parameter bebas dalam solusi,
dan perilaku hasil kali matriks. Bab ini mengembangkan hubungan tersebut
melalui ruang kolom, kernel, faktorisasi, dan eliminasi blok, kemudian
menerapkannya pada persamaan serta pangkat matriks.

Semua matriks berentri dalam $\F=\R$ atau $\C$.
Notasi $\Mat_{m\times n}(\F)$ menyatakan himpunan matriks berukuran
$m\times n$; $\Mat_n(\F)$ digunakan untuk matriks persegi.
Operasi matriks, kebebasan linear, basis, dan dimensi dianggap sudah dikenal.
Determinan belum diperlukan dalam pembuktian-pembuktian bab ini.

\section{Ruang Kolom, Ruang Baris, dan Kernel}
\label{rank:dasar}

\begin{definisi}[][Rank dan nulitas][rank:definisi]
    Untuk $A\in\Mat_{m\times n}(\F)$, definisikan
    \[
        \begin{aligned}
            \mathcal C(A)&=\{Ax:x\in\F^n\}=\operatorname{Im}A,\\
            \mathcal R(A)&=\text{ruang yang direntang oleh baris-baris }A,\\
            \ker A&=\{x\in\F^n:Ax=0\}.
        \end{aligned}
    \]
    \emph{Rank} dan \emph{nulitas} $A$ berturut-turut adalah
    $\rank A=\dim\mathcal C(A)$ dan $\nul A=\dim\ker A$.
    Matriks mempunyai rank kolom penuh jika $\rank A=n$, dan rank baris
    penuh jika $\rank A=m$.
\end{definisi}

Jadi $0\le\rank A\le\min\{m,n\}$, dan $\rank A=0$ tepat ketika $A=0$.
Rank bukan banyaknya kolom tak nol: beberapa kolom tak nol dapat
menghasilkan arah yang sama.

\begin{teorema}[][Kesamaan rank baris dan kolom][rank:baris-kolom]
    Untuk setiap matriks $A$,
    \[
        \dim\mathcal R(A)=\dim\mathcal C(A),\qquad
        \rank A^T=\rank A.
    \]
\end{teorema}
\begin{proof}
    Misalkan $c=\dim\mathcal C(A)>0$. Pilih $c$ kolom bebas dari $A$
    dan tempatkan sebagai kolom-kolom matriks $U$. Setiap kolom $A$
    merupakan kombinasi kolom $U$, sehingga $A=UV$ untuk suatu matriks
    $V$ dengan $c$ baris. Setiap baris $A$ merupakan kombinasi baris $V$;
    akibatnya $\dim\mathcal R(A)\le c$.
    Terapkan argumen yang sama pada $A^T$ untuk memperoleh ketaksamaan
    sebaliknya. Untuk $A=0$, kedua dimensi sama dengan nol.
\end{proof}

\begin{teorema}[][Rank--nulitas][rank:nulitas]
    Jika $A$ mempunyai $n$ kolom, maka
    \[
        \rank A+\nul A=n.
    \]
\end{teorema}
\begin{proof}
    Pilih basis $u_1,\ldots,u_k$ untuk $\ker A$.
    Lengkapi dengan $v_1,\ldots,v_{n-k}$ hingga menjadi basis $\F^n$.
    Vektor $Av_1,\ldots,Av_{n-k}$ merentang $\mathcal C(A)$.
    Jika $\sum_j c_jAv_j=0$, maka $\sum_j c_jv_j\in\ker A$.
    Kebebasan basis yang diperluas memaksa semua $c_j=0$.
    Jadi $\rank A=n-k$ dan $\nul A=k$.
\end{proof}

Dalam rumus ini, $n$ adalah banyaknya \emph{kolom}, yaitu dimensi
domain pemetaan $x\mapsto Ax$. Kernel adalah himpunan vektor, sedangkan
nulitas adalah dimensinya. Secara khusus, $\nul A=0$ ekuivalen dengan
$\ker A=\{0\}$ dan kebebasan semua kolom $A$.
Jika $n>m$, selalu ada solusi tak nol untuk $Ax=0$, sebab
$\nul A=n-\rank A\ge n-m>0$.

\begin{prop}[Operasi yang mempertahankan rank]
    \label{rank:invarian}
    Jika $P\in\Mat_m(\F)$ dan $Q\in\Mat_n(\F)$ invertibel, maka
    \[
        \rank(PAQ)=\rank A.
    \]
    Khususnya, operasi baris dan kolom elementer mempertahankan rank.
\end{prop}
\begin{proof}
    Karena $Q$ memetakan $\F^n$ secara bijektif ke dirinya sendiri,
    $\mathcal C(AQ)=\mathcal C(A)$. Sementara itu,
    $\mathcal C(PA)=P\mathcal C(A)$ mempunyai dimensi yang sama dengan
    $\mathcal C(A)$ karena $P$ invertibel.
\end{proof}

Pada bentuk eselon baris, banyaknya pivot sama dengan rank. Jika yang
dicari adalah basis ruang kolom matriks \emph{asal}, ambil kolom asal
yang indeksnya bersesuaian dengan pivot. Operasi baris mempertahankan
hubungan ketergantungan antarkolom, tetapi dapat mengubah ruang kolomnya.

\begin{exam}[Parameter yang mengubah rank dan nulitas]
    Tentukan rank, nulitas, dan kernel dari
    \[
        A(t)=\begin{pmatrix}1&1&5\\1&3&3\\0&-1&t\end{pmatrix},\qquad t\in\R.
    \]
\end{exam}
\begin{solution}
    Lakukan $R_2\leftarrow R_2-R_1$, kemudian
    $R_3\leftarrow R_3+\tfrac12R_2$:
    \[
        A(t)\sim
        \begin{pmatrix}1&1&5\\0&2&-2\\0&0&t-1\end{pmatrix}.
    \]
    Jika $t\ne1$, terdapat tiga pivot: $\rank A(t)=3$ dan
    $\nul A(t)=0$. Jika $t=1$, ranknya $2$, nulitasnya $1$, dan
    \[
        \ker A(1)=\operatorname{span}\{(-6,1,1)^T\}.
    \]
    Dua kolom pertama matriks asal menjadi basis ruang kolom pada $t=1$.
    Jangan membagi baris dengan $t-1$ sebelum memisahkan kasus $t=1$.
\end{solution}

\subsection{Basis kernel dan variabel bebas}
\label{rank:basis-kernel}

\begin{prop}[Membaca kernel dari bentuk eselon tereduksi]
    Misalkan $A$ mempunyai $n$ kolom dan rank $r$.
    Setelah operasi baris dan pengurutan ulang kolom, tuliskan
    \[
        EA\Pi=\begin{pmatrix}I_r&H\\0&0\end{pmatrix},
    \]
    dengan $E$ invertibel dan $\Pi$ matriks permutasi.
    Maka
    \[
        \ker A=\left\{\Pi\begin{pmatrix}-Hz\\z\end{pmatrix}:
            z\in\F^{n-r}\right\}.
    \]
    Kolom-kolom $\Pi\binom{-H}{I_{n-r}}$ membentuk basis kernel.
\end{prop}
\begin{proof}
    Gunakan perubahan peubah $x=\Pi y$, lalu pisahkan $y$ menjadi
    variabel pivot $y_1$ dan variabel bebas $y_2$.
    Persamaan $Ax=0$ ekuivalen dengan $y_1+Hy_2=0$.
    Jadi $y_2=z$ bebas dan $y_1=-Hz$.
    Pemetaan $z\mapsto\Pi(-Hz,z)^T$ injektif karena $\Pi$ invertibel
    dan komponen terakhir sebelum permutasi adalah $z$ sendiri.
    Bayangan basis standar $\F^{n-r}$ karenanya menjadi basis kernel.
\end{proof}

Jika hanya operasi baris yang dilakukan, $\ker(EA)=\ker A$.
Jika kolom diubah, koordinat solusi harus dikembalikan:
$\ker(AQ)=Q^{-1}\ker A$ untuk $Q$ invertibel.
Jumlah variabel bebas tetap sama, tetapi vektor basis kernelnya
dapat berubah dalam koordinat semula.

\subsection{Nulitas kanan dan nulitas kiri}

Nulitas $A$ biasanya berarti dimensi kernel kanan $\ker A$.
Relasi linear di antara baris-baris $A$ dicatat oleh kernel kiri,
yaitu $\ker A^T$. Untuk $A$ berukuran $m\times n$ dan rank $r$,
\[
    \nul A=n-r,\qquad \nul A^T=m-r,\qquad
    \nul A-\nul A^T=n-m.
\]
Jadi kesamaan rank $A$ dan $A^T$ tidak menjamin kesamaan nulitas
jika matriksnya persegi panjang. Di atas $\C$, kernel kiri juga
sering dinyatakan memakai $A^*=\overline A^{\,T}$ dan pasangan $y^*A=0$;
dimensinya tetap $m-r$.

\begin{exam}[Tiga variabel bebas dan satu relasi baris]
    Tentukan rank, nulitas, dan basis kernel matriks
    \[
        A=\begin{pmatrix}
            1&2&0&-1&3\\0&0&1&2&-1\\2&4&1&0&5
        \end{pmatrix}.
    \]
    Tentukan juga syarat agar $Ax=b$ konsisten.
\end{exam}
\begin{solution}
    Baris ketiga adalah dua kali baris pertama ditambah baris kedua.
    Dua baris pertama bebas, sehingga $\rank A=2$, $\nul A=5-2=3$,
    dan $\nul A^T=3-2=1$.
    Untuk $Ax=0$, ambil $x_2=s$, $x_4=t$, dan $x_5=u$.
    Maka $x_1=-2s+t-3u$ dan $x_3=-2t+u$, sehingga basis kernelnya
    \[
        \begin{pmatrix}-2\\1\\0\\0\\0\end{pmatrix},\quad
        \begin{pmatrix}1\\0\\-2\\1\\0\end{pmatrix},\quad
        \begin{pmatrix}-3\\0\\1\\0\\1\end{pmatrix}.
    \]
    Relasi baris memberi $\ker A^T=\operatorname{span}\{(-2,-1,1)^T\}$.
    Karena dua persamaan pertama dapat diselesaikan untuk setiap
    $b_1,b_2$, syarat konsistensi tepat $b_3=2b_1+b_2$.
    Jika syarat ini dipenuhi, seluruh solusi adalah
    \[
        x=\begin{pmatrix}b_1-2s+t-3u\\s\\b_2-2t+u\\t\\u\end{pmatrix},
        \qquad s,t,u\in\F.
    \]
    Tiga parameter bebas sesuai dengan nulitas $A$.
\end{solution}

Contoh ini memperlihatkan hubungan antara kernel, kebebasan kolom,
dan teorema rank--nulitas.

\section{Bentuk Normal dan Faktorisasi Rank}
\label{rank:faktorisasi}

\begin{teorema}[][Bentuk normal terhadap operasi baris dan kolom][rank:normal]
    Jika $A\in\Mat_{m\times n}(\F)$ mempunyai rank $r$, terdapat matriks
    invertibel $P,Q$ sehingga
    \[
        PAQ=\begin{pmatrix}I_r&0\\0&0\end{pmatrix}=:D_r.
    \]
    Dua matriks dengan ukuran sama mempunyai rank sama jika dan hanya
    jika dapat dihubungkan oleh perkalian matriks invertibel di kiri dan kanan.
\end{teorema}
\begin{proof}
    Reduksi baris dan permutasi kolom mengubah $A$ menjadi
    $\begin{pmatrix}I_r&H\\0&0\end{pmatrix}$.
    Kurangkan kombinasi kolom-kolom pivot dari kolom-kolom sisanya untuk
    meniadakan $H$. Semua operasi tersebut invertibel dan menghasilkan
    $D_r$. Pernyataan terakhir mengikuti bentuk normal ini dan
    \autoref{rank:invarian}.
\end{proof}

Di sini $P$ dan $Q$ boleh dipilih secara independen. Hubungan ini disebut
\emph{ekuivalensi matriks}; syaratnya berbeda dari keserupaan
$B=S^{-1}AS$, yang mewajibkan kedua faktor saling invers.

\begin{teorema}[][Faktorisasi rank penuh][rank:full-factor]
    Jika $\rank A=r>0$, maka $A$ dapat ditulis sebagai
    \[
        A=UV,\qquad U\in\Mat_{m\times r}(\F),\quad
        V\in\Mat_{r\times n}(\F),
    \]
    dengan $\rank U=\rank V=r$.
\end{teorema}
\begin{proof}
    Ambil $U$ dari $r$ kolom bebas milik $A$, lalu isi kolom ke-$j$ dari
    $V$ dengan koefisien kolom ke-$j$ dari $A$ terhadap basis tersebut.
    Pada indeks kolom yang dipilih, $V$ memuat kolom-kolom $I_r$.
    Maka $V$ mempunyai rank baris penuh dan $A=UV$.
\end{proof}

\begin{prop}[Invers satu sisi dan pencoretan faktor]
    \label{rank:inverse-sisi}
    Matriks $U$ dengan rank kolom penuh mempunyai invers kiri $L$, yaitu
    $LU=I$. Matriks $V$ dengan rank baris penuh mempunyai invers kanan
    $R$, yaitu $VR=I$. Akibatnya, untuk ukuran yang sesuai,
    \[
        \rank(UAV)=\rank A.
    \]
\end{prop}
\begin{proof}
    Dari $PUQ=\binom{I}{0}$, diperoleh
    $L=Q\begin{pmatrix}I&0\end{pmatrix}P$ dengan $LU=I$.
    Transpose memberikan pernyataan untuk invers kanan.
    Pemetaan $x\mapsto Ux$ injektif, sehingga
    $\dim U\mathcal C(A)=\dim\mathcal C(A)$.
    Karena $V$ surjektif, $\mathcal C(AV)=\mathcal C(A)$.
    Gabungkan keduanya.
\end{proof}

\begin{prop}[Submatriks invertibel terbesar]
    Rank $A$ adalah ukuran terbesar submatriks persegi invertibel
    yang dapat dipilih dari $A$, dengan nilai $0$ jika $A=0$.
\end{prop}
\begin{proof}
    Menghapus kolom tidak menambah dimensi ruang kolom; menghapus baris
    juga tidak menaikkan rank, dengan argumen yang sama pada transpose.
    Jadi submatriks persegi invertibel berukuran $s$ memaksa $s\le\rank A$.
    Sebaliknya, pilih $r=\rank A>0$ kolom bebas dan bentuk matriks $U$.
    Rank baris $U$ juga $r$, sehingga dapat dipilih $r$ baris bebasnya.
    Irisan baris dan kolom yang terpilih adalah matriks persegi dengan
    rank penuh, sehingga invertibel oleh rank--nulitas.
\end{proof}

\subsection{Membaca rank dari pola entri}

Jika $U=[u_1\ \cdots\ u_r]$ dan baris-baris $V$ adalah
$v_1^T,\ldots,v_r^T$, maka
\[
    A=\sum_{j=1}^r u_jv_j^T.
\]
Setiap $u_jv_j^T$ mempunyai rank satu jika kedua vektornya tak nol.
Sebaliknya, jika $A$ merupakan jumlah $s$ matriks rank satu, ruang kolomnya
direntang oleh paling banyak $s$ vektor. Jadi rank $A$ adalah banyaknya
suku rank satu paling sedikit yang diperlukan untuk menyatakan $A$.

\begin{exam}[Matriks besar dengan rank dua]
    Untuk $n\ge2$, tentukan rank matriks $A=(i+j-1)_{i,j=1}^n$.
\end{exam}
\begin{solution}
    Definisikan $u=(1,2,\ldots,n)^T$ dan $e=(1,\ldots,1)^T$.
    Kolom ke-$j$ dari $A$ adalah $u+(j-1)e$, sehingga
    $\mathcal C(A)\subseteq\operatorname{span}\{u,e\}$.
    Dua kolom pertama bebas karena $u$ dan $e$ bebas.
    Maka $\rank A=2$, berapa pun besarnya $n$.
\end{solution}

\begin{exam}[Pola kuadratik tanpa menghitung determinan]
    Untuk $n\ge3$, matriks $B=(i^2+ij+j^2)_{i,j=1}^n$ mempunyai rank $3$.
\end{exam}
\begin{solution}
    Faktorkan
    \[
        B=\begin{pmatrix}
            1&1&1\\1&2&2^2\\\vdots&\vdots&\vdots\\1&n&n^2
        \end{pmatrix}
        \begin{pmatrix}
            1^2&2^2&\cdots&n^2\\1&2&\cdots&n\\1&1&\cdots&1
        \end{pmatrix}=UV.
    \]
    Kolom $U$ bebas: polinom berderajat paling tinggi dua yang nol di
    $1,2,3$ harus identik nol. Alasan sama membuktikan baris $V$ bebas.
    Dengan \autoref{rank:inverse-sisi}, $\rank(UV)=3$.
    Faktorisasi memberi batas atas; kebebasan kedua faktor memastikan
    bahwa batas tersebut benar-benar tercapai.
\end{solution}

\section{Rank dan Nulitas Hasil Kali}
\label{rank:produk}

\begin{teorema}[][Kehilangan dimensi pada hasil kali][rank:produk-tepat]
    Untuk $A\in\Mat_{m\times n}(\F)$ dan $B\in\Mat_{n\times p}(\F)$,
    \[
        \rank(AB)=\rank B-\dim\bigl(\mathcal C(B)\cap\ker A\bigr).
    \]
\end{teorema}
\begin{proof}
    Batasi pemetaan $x\mapsto Ax$ pada domain $\mathcal C(B)$.
    Citranya adalah $\mathcal C(AB)$, sedangkan kernelnya
    $\mathcal C(B)\cap\ker A$. Terapkan rank--nulitas pada pemetaan ini.
\end{proof}

\begin{prop}[Nulitas hasil kali]
    \label{rank:nulitas-produk}
    Untuk $A$ berukuran $m\times n$ dan $B$ berukuran $n\times p$,
    \[
        \nul(AB)=\nul B+\dim\bigl(\mathcal C(B)\cap\ker A\bigr).
    \]
    Akibatnya,
    \[
        \nul B\le\nul(AB)\le\min\{p,\nul A+\nul B\}.
    \]
    Jika $B$ mempunyai rank baris penuh, maka
    $\nul(AB)=\nul A+\nul B$.
\end{prop}
\begin{proof}
    Matriks $B$ dan $AB$ sama-sama mempunyai $p$ kolom.
    Kurangkan rumus pada \autoref{rank:produk-tepat} dari $p$.
    Untuk batasnya, gunakan
    $0\le\dim(\mathcal C(B)\cap\ker A)\le\nul A$.
    Jika $B$ mempunyai rank baris penuh, $\mathcal C(B)=\F^n$,
    sehingga irisannya dengan $\ker A$ adalah seluruh $\ker A$.
\end{proof}

\begin{exam}[Kernel tambahan pada hasil kali]
    Misalkan
    \[
        A=\begin{pmatrix}1&1&0\\0&0&0\end{pmatrix},\qquad
        B=\begin{pmatrix}1&0&1&0\\0&1&1&0\\0&0&0&0\end{pmatrix}.
    \]
    Tentukan nulitas $A$, $B$, dan $AB$ melalui ruang kolom dan kernel.
\end{exam}
\begin{solution}
    Rank $A$ adalah $1$ dan rank $B$ adalah $2$, sehingga
    $\nul A=3-1=2$ dan $\nul B=4-2=2$.
    Selanjutnya,
    \[
        \mathcal C(B)=\operatorname{span}\{e_1,e_2\},\qquad
        \ker A=\operatorname{span}\{e_1-e_2,e_3\}.
    \]
    Irisannya berdimensi satu. Jadi $\nul(AB)=2+1=3$ dan $\rank(AB)=1$.
    Selain arah dalam $\ker B$, ada satu arah bebas tambahan yang
    keluaran $B$-nya dimatikan oleh $A$.
    Arah $e_3\in\ker A$ tidak menyumbang kehilangan tambahan karena
    tidak termasuk ruang kolom $B$.
\end{solution}

\ide{Hasil kali $AB$ kehilangan rank ketika sebagian keluaran $B$
    masuk ke kernel $A$. Rank kedua faktor saja belum menentukan
    besar kehilangan tersebut.}

\begin{teorema}[][Ketaksamaan Sylvester][rank:sylvester]
    Dengan ukuran seperti di atas,
    \[
        \max\{0,\rank A+\rank B-n\}
        \le\rank(AB)\le\min\{\rank A,\rank B\}.
    \]
    Pada batas bawah, $n$ adalah dimensi \emph{ruang perantara}:
    banyak kolom $A$, sekaligus banyak baris $B$.
\end{teorema}
\begin{proof}
    Gunakan $\mathcal C(AB)\subseteq\mathcal C(A)$ dan
    \autoref{rank:produk-tepat}. Untuk batas bawah,
    \[
        \dim(\mathcal C(B)\cap\ker A)\le\dim\ker A=n-\rank A.
    \]
\end{proof}

\begin{prop}[Kapan batas hasil kali tercapai?]
    \label{rank:kesamaan-produk}
    Untuk ukuran yang sesuai,
    \begin{align*}
        \rank(AB)=\rank B
        &\iff \mathcal C(B)\cap\ker A=\{0\},\\
        \rank(AB)=\rank A
        &\iff \mathcal C(B)+\ker A=\F^n,\\
        \rank(AB)=\rank A+\rank B-n
        &\iff \ker A\subseteq\mathcal C(B).
    \end{align*}
\end{prop}
\begin{proof}
    Pernyataan pertama dan ketiga langsung mengikuti rumus rank hasil
    kali. Untuk yang kedua, $A\mathcal C(B)=\mathcal C(A)$ tepat ketika
    setiap $x\in\F^n$ dapat ditulis $x=y+z$ dengan
    $y\in\mathcal C(B)$ dan $z\in\ker A$.
\end{proof}

Khususnya, $AB=0$ ekuivalen dengan $\mathcal C(B)\subseteq\ker A$,
sehingga $\rank A+\rank B\le n$.
Ketaksamaan rank tersebut saja tidak cukup untuk menyimpulkan $AB=0$:
ambil $A=B=\diag(1,0)$.

\begin{exam}[Seluruh kemungkinan rank hasil kali]
    Misalkan $A,B\in\Mat_8(\R)$, $\rank A=5$, dan $\rank B=6$.
    Tentukan semua nilai yang mungkin bagi $\rank(AB)$.
\end{exam}
\begin{solution}
    Sylvester memberi $3\le\rank(AB)\le5$.
    Untuk membuktikan bahwa $3,4,5$ semuanya mungkin, tetapkan
    $A=\diag(I_5,0_3)$. Untuk setiap $k\in\{3,4,5\}$, pilih ruang
    kolom $B$ yang direntang oleh $k$ vektor koordinat dari
    $e_1,\ldots,e_5$ dan $6-k$ vektor dari $e_6,e_7,e_8$.
    Jadikan keenam vektor tersebut kolom pertama $B$, lalu tambahkan
    dua kolom nol. Maka $\rank B=6$ dan tepat $k$ arah bertahan pada $AB$.
    Secara umum, cara ini merealisasikan setiap bilangan bulat antara
    $\max\{0,a+b-n\}$ dan $\min\{a,b\}$ untuk dua matriks persegi
    berukuran $n$ dengan rank $a,b$.
\end{solution}

\begin{teorema}[][Ketaksamaan Frobenius][rank:frobenius]
    Jika $AB$ dan $BC$ terdefinisi, maka
    \[
        \rank(ABC)\ge\rank(AB)+\rank(BC)-\rank B.
    \]
\end{teorema}
\begin{proof}
    Karena $\mathcal C(BC)\subseteq\mathcal C(B)$,
    \begin{align*}
        \rank(ABC)
        &=\rank(BC)-\dim(\mathcal C(BC)\cap\ker A)\\
        &\ge\rank(BC)-\dim(\mathcal C(B)\cap\ker A)\\
        &=\rank(BC)+\rank(AB)-\rank B.
    \end{align*}
\end{proof}

Batas Frobenius mempertahankan informasi tentang faktor tengah $B$.
Jika hanya ukuran ruang perantara yang dipakai, informasi itu dapat
hilang ketika Sylvester diterapkan secara berulang.

\section{Rank Jumlah dan Perubahan Rank Satu}
\label{rank:jumlah}

\begin{prop}[Matriks gabungan]
    \label{rank:gabungan}
    Untuk $A,B$ dengan jumlah baris sama,
    \[
        \rank[A\ B]=\rank A+\rank B
        -\dim(\mathcal C(A)\cap\mathcal C(B)).
    \]
    Untuk $A,B$ dengan jumlah kolom sama, berlaku
    \[
        \ker\begin{pmatrix}A\\B\end{pmatrix}=\ker A\cap\ker B.
    \]
\end{prop}
\begin{proof}
    Ruang kolom $[A\ B]$ adalah $\mathcal C(A)+\mathcal C(B)$.
    Gunakan identitas
    $\dim(U+V)=\dim U+\dim V-\dim(U\cap V)$, yang diperoleh dengan
    memperluas basis $U\cap V$ menjadi basis masing-masing ruang.
    Pernyataan kernel mengikuti kesetaraan
    $\binom{A}{B}x=0\iff Ax=Bx=0$.
\end{proof}

\begin{teorema}[][Batas rank jumlah][rank:batas-jumlah]
    Jika $A,B\in\Mat_{m\times n}(\F)$, maka
    \[
        |\rank A-\rank B|\le\rank(A+B)
        \le\min\{m,n,\rank A+\rank B\}.
    \]
    Ketaksamaan yang sama berlaku untuk $A-B$.
\end{teorema}
\begin{proof}
    Inklusi $\mathcal C(A+B)\subseteq\mathcal C(A)+\mathcal C(B)$
    memberi batas atas. Terapkan batas ini pada
    $A=(A+B)-B$ dan $B=(A+B)-A$ untuk memperoleh batas bawah.
\end{proof}

\begin{teorema}[][Kapan rank bersifat aditif?][rank:aditif]
    Untuk dua matriks berukuran sama,
    \[
        \rank(A+B)=\rank A+\rank B
    \]
    jika dan hanya jika kedua syarat berikut dipenuhi:
    \[
        \mathcal C(A)\cap\mathcal C(B)=\{0\},\qquad
        \mathcal R(A)\cap\mathcal R(B)=\{0\}.
    \]
\end{teorema}
\begin{proof}
    Dari $A+B=[A\ B]\binom{I}{I}$ dan
    $A+B=[I\ I]\binom{A}{B}$, rank jumlah tidak melebihi rank kedua
    matriks gabungan. Jika ranknya $\rank A+\rank B$, rumus dimensi
    ruang kolom dan versi transposenya memaksa kedua irisan menjadi nol.

    Sebaliknya, irisan ruang baris yang nol memberi
    $\rank\binom{A}{B}=\rank A+\rank B$.
    Pemetaan $(u,v)\mapsto u+v$ injektif pada
    $\mathcal C(A)\times\mathcal C(B)$ karena irisan ruang kolom nol.
    Maka penerapan pemetaan tersebut pada citra $\binom{A}{B}$ tidak
    menurunkan dimensi, sehingga rank $A+B$ sama dengan jumlah ranknya.
\end{proof}

Satu syarat saja tidak cukup. Misalnya,
\[
    A=\begin{pmatrix}1&0\\0&0\end{pmatrix},\qquad
    B=\begin{pmatrix}0&0\\1&0\end{pmatrix}
\]
mempunyai ruang kolom yang beririsan nol, tetapi
$\rank(A+B)=1<2=\rank A+\rank B$ karena ruang barisnya berimpit.

\subsection{Mengganggu matriks dengan satu hasil kali luar}

Karena $\rank(uv^T)\le1$, selalu berlaku
\[
    \bigl|\rank(A+uv^T)-\rank A\bigr|\le1.
\]
Untuk matriks invertibel, kasus penurunan rank dapat ditentukan tepat.

\begin{prop}[Perubahan rank satu pada matriks invertibel]
    \label{rank:update}
    Jika $A\in\Mat_n(\F)$ invertibel dan $u,v\in\F^n$, maka
    \[
        \rank(A+uv^T)=
        \begin{cases}
            n,&1+v^TA^{-1}u\ne0,\\
            n-1,&1+v^TA^{-1}u=0.
        \end{cases}
    \]
    Pada kasus kedua, $\ker(A+uv^T)=\operatorname{span}\{A^{-1}u\}$.
\end{prop}
\begin{proof}
    Tuliskan $w=A^{-1}u$. Persamaan $(A+uv^T)x=0$ ekuivalen dengan
    $x=-w(v^Tx)$. Jadi semua vektor kernel berada pada
    $\operatorname{span}\{w\}$. Mengalikan dengan $v^T$ memberi
    $(1+v^Tw)v^Tx=0$.
    Jika koefisiennya tak nol, $x=0$.
    Jika koefisiennya nol, $w\ne0$ dan $(I+wv^T)w=0$;
    kernel berdimensi satu. Gunakan rank--nulitas.
\end{proof}

\begin{exam}[Diagonal konstan dan satu matriks rank satu]
    Misalkan $J_n$ adalah matriks yang semua entrinya $1$, $n\ge2$.
    Tentukan rank $tI_n+J_n$ untuk setiap $t\in\F$.
\end{exam}
\begin{solution}
    Tuliskan $J_n=ee^T$ dengan $e=(1,\ldots,1)^T$.
    Jika $t\ne0$, terapkan \autoref{rank:update} pada $A=tI_n$;
    skalar penentunya adalah $1+n/t$.
    Kasus $t=0$ dihitung langsung. Hasilnya
    \[
        \rank(tI_n+J_n)=
        \begin{cases}
            1,&t=0,\\
            n-1,&t=-n,\\
            n,&t\notin\{0,-n\}.
        \end{cases}
    \]
    Pada $t=-n$ kernelnya $\operatorname{span}\{e\}$; pada $t=0$
    kernelnya $\{x:e^Tx=0\}$.
    Ini juga menyelesaikan matriks dengan entri diagonal $a$ dan entri
    luar diagonal $b$: matriksnya adalah $(a-b)I_n+bJ_n$.
\end{solution}

\section{Eliminasi Blok dan Persamaan Matriks}
\label{rank:blok}

Operasi baris blok dilakukan dengan mengalikan blok dari kiri;
operasi kolom blok mengalikan dari kanan. Urutan perkalian harus dijaga
karena dua matriks umumnya tidak komutatif.

\begin{teorema}[][Rumus rank komplemen Schur][rank:schur]
    Misalkan
    $M=\begin{pmatrix}P&Q\\R&S\end{pmatrix}$, dengan $P$ persegi
    berukuran $k$ dan invertibel. Maka
    \[
        \rank M=k+\rank(S-RP^{-1}Q).
    \]
\end{teorema}
\begin{proof}
    Identitas
    \[
        \begin{pmatrix}I&0\\-RP^{-1}&I\end{pmatrix}
        \begin{pmatrix}P&Q\\R&S\end{pmatrix}
        \begin{pmatrix}I&-P^{-1}Q\\0&I\end{pmatrix}
        =\begin{pmatrix}P&0\\0&S-RP^{-1}Q\end{pmatrix}
    \]
    menggunakan dua faktor invertibel di kiri dan kanan. Rank matriks
    diagonal blok adalah jumlah rank blok-bloknya karena ruang kolomnya
    menempati kelompok koordinat yang terpisah.
\end{proof}

\begin{prop}[Segitiga blok]
    \label{rank:segitiga-blok}
    Untuk blok-blok berukuran sesuai,
    \[
        \rank\begin{pmatrix}A&X\\0&B\end{pmatrix}
        \ge\rank A+\rank B.
    \]
    Kesamaan selalu berlaku jika $X=0$, tetapi belum tentu jika $X\ne0$.
\end{prop}
\begin{proof}
    Pilih $\rank A$ kolom bebas dari blok kolom pertama dan
    $\rank B$ kolom dari blok kolom kedua yang bagian bawahnya bebas.
    Gabungannya bebas: komponen bawah terlebih dahulu memaksa koefisien
    kelompok kedua nol, lalu komponen atas memaksa sisanya nol.
    Untuk melihat kemungkinan ketaksamaan ketat, ambil
    $A=B=(0)$ dan $X=(1)$.
\end{proof}

\begin{exam}[Dua hasil kali dengan ukuran berbeda]
    Jika $A$ berukuran $m\times n$ dan $B$ berukuran $n\times m$, buktikan
    \[
        m-\rank(I_m-AB)=n-\rank(I_n-BA).
    \]
\end{exam}
\begin{solution}
    Hitung rank matriks yang sama dengan dua pilihan pivot blok:
    \[
        \rank\begin{pmatrix}I_m&A\\B&I_n\end{pmatrix}
        =m+\rank(I_n-BA)=n+\rank(I_m-AB).
    \]
    Jadi kedua matriks mempunyai nulitas sama, walaupun ukurannya
    berbeda. Secara langsung, $x\mapsto Bx$ memberi isomorfisme dari
    $\ker(I_m-AB)$ ke $\ker(I_n-BA)$ dengan invers $y\mapsto Ay$.
    Pernyataan ini tidak berarti $\rank(AB)=\rank(BA)$.
\end{solution}

\subsection{Konsistensi dan banyaknya parameter bebas}

\begin{prop}[Persamaan dengan satu faktor tetap]
    \label{rank:ax-b}
    Untuk $A\in\Mat_{m\times n}(\F)$ dan $C\in\Mat_{m\times q}(\F)$,
    persamaan $AX=C$ mempunyai solusi jika dan hanya jika
    \[
        \mathcal C(C)\subseteq\mathcal C(A)
        \iff \rank[A\ C]=\rank A.
    \]
    Jika konsisten, himpunan solusinya berbentuk $X_0+\mathcal N$, dengan
    $\mathcal N=\{Z:AZ=0\}$ berdimensi $q\,\nul A=q(n-\rank A)$.
\end{prop}
\begin{proof}
    Persamaan ini setara dengan $q$ persamaan kolom $Ax_j=c_j$.
    Setiap kolom $C$ harus berada pada $\mathcal C(A)$.
    Jika $X_0$ adalah satu solusi, seluruh solusi adalah $X_0+Z$,
    dengan setiap kolom $Z$ dipilih secara bebas dari $\ker A$.
    Tiap kolom menyumbang $n-\rank A$ parameter bebas.
\end{proof}

Untuk $q=1$, diperoleh kriteria konsistensi sistem linear $Ax=c$.
Jika konsisten, solusi tunggal tepat ketika $\nul A=0$, atau ekuivalen
$\rank A=n$.
Untuk $C\ne0$, himpunan solusi umumnya merupakan ruang afin
(translasi suatu subruang), bukan subruang yang melalui nol.

\begin{teorema}[][Persamaan dengan dua faktor tetap][rank:axb]
    Ambil $A\in\Mat_{m\times n}(\F)$,
    $B\in\Mat_{p\times q}(\F)$, dan $C\in\Mat_{m\times q}(\F)$.
    Persamaan $AXB=C$, dengan peubah $X\in\Mat_{n\times p}(\F)$,
    konsisten jika dan hanya jika
    \[
        \mathcal C(C)\subseteq\mathcal C(A),\qquad
        \mathcal R(C)\subseteq\mathcal R(B).
    \]
    Jika konsisten, banyak parameter bebasnya adalah
    \[
        np-(\rank A)(\rank B).
    \]
\end{teorema}
\begin{proof}
    Syarat inklusi perlu karena perkalian di kanan menghasilkan
    kombinasi kolom, sedangkan perkalian di kiri menghasilkan kombinasi baris.
    Untuk membuktikan kecukupan dan menghitung dimensinya, pilih
    \[
        PAQ=D_a,\qquad RBS=D_b,\qquad a=\rank A,\quad b=\rank B,
    \]
    dalam bentuk normal \autoref{rank:normal}.
    Dengan $Y=Q^{-1}XR^{-1}$ dan $C'=PCS$, persamaannya menjadi
    \[
        D_aYD_b=C'.
    \]
    Kedua syarat inklusi tepat menyatakan bahwa entri $C'$ di luar
    $a$ baris pertama dan $b$ kolom pertama semuanya nol.
    Perkalian $D_aYD_b$ mengambil blok $a\times b$ di pojok kiri atas
    $Y$ dan mengisi bagian lainnya dengan nol.
    Jadi $ab$ entri $Y$ ditentukan oleh $C'$, sedangkan $np-ab$ entri
    sisanya bebas. Perubahan peubah antara $X$ dan $Y$ bijektif.
\end{proof}

Hasil ini merupakan penerapan bentuk normal rank dan kriteria
faktorisasi ruang kolom dan baris.
Kondisi di atas juga dapat diperiksa melalui
$\rank[A\ C]=\rank A$ dan
$\rank\binom{B}{C}=\rank B$.

Pandang $\mathcal L(X)=AXB$ sebagai pemetaan linear dari
$\Mat_{n\times p}(\F)$ ke $\Mat_{m\times q}(\F)$.
Domainnya berdimensi $np$, dan pembuktian di atas menunjukkan
\[
    \rank\mathcal L=(\rank A)(\rank B),\qquad
    \nul\mathcal L=np-(\rank A)(\rank B).
\]
Di sini rank dan nulitas pemetaan berarti dimensi citra dan kernelnya.
Nulitas tersebut adalah dimensi ruang matriks $Z$ yang memenuhi $AZB=0$,
sekaligus banyaknya parameter bebas pada setiap persamaan $AXB=C$
yang konsisten.

\begin{exam}[Sembilan peubah, empat kendala bebas]
    Misalkan
    \[
        A=\begin{pmatrix}1&2&0\\0&0&1\end{pmatrix},\qquad
        B=\begin{pmatrix}1&0\\0&1\\1&1\end{pmatrix},\qquad
        C\in\Mat_2(\R).
    \]
    Tentukan apakah $AXB=C$ selalu dapat diselesaikan dan berapa banyak
    parameter bebas dalam solusinya.
\end{exam}
\begin{solution}
    Matriks $A$ mempunyai rank baris penuh $2$, dan $B$ mempunyai
    rank kolom penuh $2$. Maka kedua syarat inklusi selalu dipenuhi,
    dan banyak parameter bebasnya $3\cdot3-2\cdot2=5$.
    Secara konstruktif, gunakan
    \[
        R=\begin{pmatrix}1&0\\0&0\\0&1\end{pmatrix},\qquad
        L=\begin{pmatrix}1&0&0\\0&1&0\end{pmatrix}.
    \]
    Karena $AR=I_2$ dan $LB=I_2$, matriks $X_0=RCL$ adalah satu solusi.
    Semua solusi diperoleh dengan menambahkan $Z$ yang memenuhi $AZB=0$.
\end{solution}

\section{Transpose Konjugat dan Pangkat Matriks}
\label{rank:pangkat}

\subsection{Menggunakan kepositifan untuk menemukan kernel}

Untuk matriks kompleks, tuliskan $A^*=\overline A^{\,T}$.
Jika $A$ real, $A^*=A^T$.

\begin{teorema}[][Rank matriks Gram][rank:gram]
    Untuk setiap matriks kompleks $A$,
    \[
        \ker(A^*A)=\ker A,\qquad
        \rank(A^*A)=\rank A=\rank(AA^*).
    \]
\end{teorema}
\begin{proof}
    Jelas $Ax=0$ mengakibatkan $A^*Ax=0$.
    Sebaliknya, jika $A^*Ax=0$, maka
    \[
        0=x^*A^*Ax=(Ax)^*(Ax)=\sum_i |(Ax)_i|^2,
    \]
    sehingga $Ax=0$.
    Rank--nulitas memberi kesamaan rank pertama.
    Konjugasi mempertahankan kebebasan linear, sehingga
    $\rank A^*=\rank A$. Terapkan hasil pertama pada $A^*$ untuk
    mendapatkan $\rank(AA^*)=\rank A$.
\end{proof}

\begin{remark}
    Di atas $\C$, jangan mengganti $A^*$ dengan $A^T$ secara sembarang.
    Untuk $A=(1,i)^T$, diperoleh $A^TA=0$, padahal $\rank A=1$.
    Sebaliknya, $A^*A=(2)$ mempunyai rank satu sebagaimana mestinya.
\end{remark}

Jika $A$ berukuran $m\times n$ dan mempunyai rank $r$, teorema Gram memberi
\[
    \nul(A^*A)=n-r=\nul A,\qquad \nul(AA^*)=m-r.
\]
Kedua matriks Gram mempunyai rank sama, sedangkan nulitasnya mengikuti
ukuran masing-masing matriks.

\subsection{Barisan rank, nulitas, dan stabilisasi kernel}

\begin{teorema}[][Penurunan rank pangkat][rank:deret-pangkat]
    Misalkan $A\in\Mat_n(\F)$ dan $r_k=\rank(A^k)$, dengan $r_0=n$.
    Maka
    \[
        n=r_0\ge r_1\ge r_2\ge\cdots\ge0,
        \qquad r_{k-1}-r_k\ge r_k-r_{k+1}\quad(k\ge1).
    \]
    Jika $r_s=r_{s+1}$ untuk suatu $s\ge0$, semua rank berikutnya
    sama dengan $r_s$.
\end{teorema}
\begin{proof}
    Inklusi
    \[
        \ker A^k\subseteq\ker A^{k+1},\qquad
        \mathcal C(A^{k+1})\subseteq\mathcal C(A^k)
    \]
    memberi penurunan rank. Terapkan Frobenius pada
    $A\cdot A^{k-1}\cdot A$:
    \[
        r_{k+1}\ge2r_k-r_{k-1}.
    \]
    Ini tepat ketaksamaan untuk selisih rank.
    Selisih-selisih tersebut nonnegatif dan tidak meningkat, sehingga
    sesudah satu selisih nol, semua selisih berikutnya juga nol.
\end{proof}

Jadi bukan setiap barisan bilangan bulat menurun dapat menjadi barisan
rank pangkat. Misalnya, untuk matriks berukuran $8$ dengan rank $6$,
harus berlaku $\rank A^2\ge4$; nilai $3$ mustahil.
Stabilisasi pasti terjadi paling lambat pada $A^n$, sebab rank tidak
dapat turun ketat lebih dari $n$ kali.

Dalam bahasa nulitas, tuliskan $d_k=\nul(A^k)=n-r_k$, dengan $d_0=0$.
Maka
\[
    0=d_0\le d_1\le d_2\le\cdots\le n,\qquad
    d_{k+1}-d_k\le d_k-d_{k-1}.
\]
Kernel bertambah besar, tetapi pertambahan dimensinya tidak meningkat.
Karena $\ker A^k\subseteq\ker A^{k+1}$, kesamaan dua nulitas berurutan
berarti kedua kernel sudah sama; setelah itu seluruh kernel berikutnya
juga sama. Selain itu, nulitas hasil kali memberi
\[
    \nul A\le\nul(A^k)\le\min\{n,k\,\nul A\}\qquad(k\ge1).
\]
Batas kanan diperoleh dengan menerapkan
\autoref{rank:nulitas-produk} secara berulang.

\begin{prop}[Rank tidak turun setelah dikuadratkan]
    \label{rank:stabil-awal}
    Untuk $A\in\Mat_n(\F)$, pernyataan berikut ekuivalen:
    \begin{enumerate}[label=(\roman*),nosep]
        \item $\rank A^2=\rank A$;
        \item $\mathcal C(A)\cap\ker A=\{0\}$;
        \item $\F^n=\mathcal C(A)\oplus\ker A$;
        \item dalam suatu basis, matriks $A$ berbentuk
            $\begin{pmatrix}D&0\\0&0\end{pmatrix}$ dengan $D$ invertibel
            berukuran $\rank A$.
    \end{enumerate}
    Untuk $A=0$, blok $D$ kosong dan pernyataan terakhir diartikan
    sebagai matriks nol.
\end{prop}
\begin{proof}
    Ekuivalensi (i)--(ii) berasal dari rumus rank hasil kali.
    Karena dimensi ruang kolom dan kernel berjumlah $n$, (ii) ekuivalen
    dengan (iii). Pada (iii), kedua subruang invarian terhadap $A$.
    Pembatasan $A$ pada $\mathcal C(A)$ injektif, sehingga invertibel
    pada ruang berdimensi hingga itu; pada $\ker A$, pemetaannya nol.
    Gabungan basis kedua subruang memberikan (iv).
    Sebaliknya, bentuk (iv) langsung memberikan (i).
\end{proof}

\begin{prop}[Batas rank dan nulitas matriks nilpoten]
    \label{rank:nilpoten}
    Jika $A^k=0$ untuk suatu bilangan bulat $k\ge1$, maka
    \[
        \rank A\le n-\left\lceil\frac nk\right\rceil,
        \qquad \nul A\ge\left\lceil\frac nk\right\rceil.
    \]
    Khususnya, $A^2=0$ mengakibatkan
    $\rank A\le\lfloor n/2\rfloor$ dan $\nul A\ge\lceil n/2\rceil$.
\end{prop}
\begin{proof}
    Sylvester berulang memberi
    $\rank A^k\ge k\rank A-(k-1)n$.
    Karena ruas kiri nol, $\rank A\le n-n/k$.
    Rank bilangan bulat, sehingga batasnya
    $\lfloor n-n/k\rfloor=n-\lceil n/k\rceil$.
    Untuk $k=2$, bukti langsungnya adalah
    $\mathcal C(A)\subseteq\ker A$.
\end{proof}

\begin{exam}[Batas yang dapat dicapai]
    Berapakah rank terbesar dan nulitas terkecil matriks $10\times10$
    yang memenuhi $A^3=0$?
\end{exam}
\begin{solution}
    Batas di atas memberi $\rank A\le10-\lceil10/3\rceil=6$.
    Definisikan matriks pergeseran
    \[
        N_3=\begin{pmatrix}0&1&0\\0&0&1\\0&0&0\end{pmatrix}.
    \]
    Matriks $A=\diag(N_3,N_3,N_3,0)$ berukuran $10$, memenuhi $A^3=0$,
    dan mempunyai rank $2+2+2=6$. Jadi rank terbesar adalah $6$ dan
    nulitas terkecil adalah $10-6=4$.
    Untuk ukuran umum, rantai pergeseran dengan panjang paling tinggi
    $k$ juga memperlihatkan ketajaman batas pada \autoref{rank:nilpoten}.
\end{solution}

\section{Identitas Matriks, Rank, dan Nulitas}
\label{rank:identitas}

\begin{teorema}[][Idempoten dan pasangan komplemennya][rank:idempoten]
    Jika $E^2=E$, maka
    \[
        \mathcal C(E)=\ker(I-E),\qquad
        \ker E=\mathcal C(I-E),\qquad
        \rank E+\rank(I-E)=n.
    \]
    Sebaliknya, jika $A+B=I_n$ dan $\rank A+\rank B=n$, maka
    $A^2=A$, $B^2=B$, dan $AB=BA=0$.
\end{teorema}
\begin{proof}
    Jika $y=Ex$, maka $Ey=E^2x=y$. Sebaliknya, $Ey=y$ berarti
    $y\in\mathcal C(E)$. Ini membuktikan kesamaan pertama.
    Kesamaan kedua mengikuti $E(I-E)=0$ dan
    $x\in\ker E\Rightarrow x=(I-E)x$.
    Rank--nulitas memberi jumlah rank yang dinyatakan.

    Untuk arah sebaliknya, $A+B=I$ memberi
    $\mathcal C(A)+\mathcal C(B)=\F^n$. Syarat jumlah dimensi
    memaksa irisannya nol. Karena $B=I-A$, berlaku $AB=BA$;
    setiap kolom $AB$ berada pada kedua ruang kolom, sehingga $AB=BA=0$.
    Maka $A=A(A+B)=A^2$ dan $B=(A+B)B=B^2$.
\end{proof}

\begin{prop}[Identitas kuadratik dengan dua akar berbeda]
    \label{rank:kuadratik}
    Jika $(A-aI)(A-bI)=0$, dengan $a\ne b$, maka
    \[
        \rank(A-aI)+\rank(A-bI)=n.
    \]
    Khususnya, jika $A^2=I$, maka
    $\rank(A-I)+\rank(A+I)=n$.
\end{prop}
\begin{proof}
    Definisikan
    \[
        E=\frac{A-bI}{a-b},\qquad F=\frac{aI-A}{a-b}.
    \]
    Kedua matriks berjumlah $I$ dan memenuhi $EF=FE=0$.
    Akibatnya $E^2=E$ dan $F=I-E$. Terapkan
    \autoref{rank:idempoten}, lalu gunakan bahwa perkalian dengan
    skalar tak nol mempertahankan rank.
\end{proof}

Sebagai contoh, jika $A\in\Mat_3(\F)$ memenuhi $A^2=I$ dan
$A\ne\pm I$, kedua rank $A-I$ dan $A+I$ positif serta berjumlah $3$.
Jadi ranknya $1$ dan $2$, dalam suatu urutan.

\begin{prop}[Identitas rank untuk $A-ABA$]
    \label{rank:aba}
    Jika $A\in\Mat_{m\times n}(\F)$ dan
    $B\in\Mat_{n\times m}(\F)$, maka
    \[
        \rank(A-ABA)=\rank A+\rank(I_n-BA)-n.
    \]
\end{prop}
\begin{proof}
    Untuk setiap $x\in\ker A$, berlaku $(I_n-BA)x=x$.
    Maka $\ker A\subseteq\mathcal C(I_n-BA)$, sehingga irisan kedua
    ruang itu tepat $\ker A$. Terapkan rumus rank hasil kali pada
    $A(I_n-BA)$:
    \[
        \rank(A-ABA)=\rank(I_n-BA)-\dim\ker A.
    \]
\end{proof}

Untuk idempoten $E$, hubungan ruang kolom dan kernel di atas juga memberi
$\nul E=\rank(I-E)$ dan $\nul E+\nul(I-E)=n$.
Untuk $(A-aI)(A-bI)=0$ dengan $a\ne b$, rank--nulitas memberikan
\[
    \nul(A-aI)+\nul(A-bI)=n.
\]
Jadi solusi kedua sistem homogen itu bersama-sama menyediakan $n$
arah bebas: kernelnya beririsan nol dan berjumlah langsung menjadi $\F^n$.
Irisannya nol karena $Ax=ax$ dan $Ax=bx$ memaksa $(a-b)x=0$.

\subsection{Memilih pendekatan pada soal rank dan nulitas}

\begin{enumerate}[label=\arabic*.,leftmargin=*,itemsep=3pt]
    \item Untuk pola entri, cari faktorisasi $UV$ atau tulis kolom
        sebagai kombinasi beberapa vektor tetap. Buktikan juga batas
        bawahnya; faktorisasi saja baru memberi batas atas.
    \item Untuk nulitas, tentukan banyak kolom lebih dahulu. Jika kernel
        mudah dideskripsikan, hitung dimensinya langsung; jika rank lebih
        mudah ditemukan, gunakan $\nul A=n-\rank A$.
    \item Untuk hasil kali, selidiki $\mathcal C(B)\cap\ker A$;
        gunakan Sylvester jika hanya rank kedua faktor yang diketahui.
    \item Untuk jumlah, periksa ruang kolom dan ruang baris sekaligus.
        Rank umumnya tidak aditif.
    \item Untuk blok, cari blok invertibel yang bisa dijadikan pivot.
        Untuk parameter, pisahkan dahulu nilai yang membuat pivot gagal.
    \item Untuk identitas seperti $A^2=0$ atau $A^2=A$, ubah identitas
        menjadi inklusi atau dekomposisi ruang kolom dan kernel.
    \item Untuk rank pangkat, periksa apakah penurunan rank mengecil
        dan apakah kernel sudah stabil.
\end{enumerate}

Hubungan dengan determinan akan dibahas pada \autoref{ch:determinan}.
Rank penuh matriks persegi berarti matriks tersebut invertibel;
nanti kondisi yang sama dapat dikenali melalui determinan tak nol.

\section{Latihan Pilihan}
\label{rank:latihan}

Soal berikut menggabungkan perhitungan struktural dan pembuktian.

\begin{latihan}\label{rank:lat-siklus}
    Misalkan $e_1,\ldots,e_n$ basis standar $\F^n$, $n\ge2$.
    Matriks $C_t$ mempunyai kolom
    \[
        e_1-e_2,\ e_2-e_3,\ldots,e_{n-1}-e_n,\ e_n-te_1.
    \]
    Tentukan rank dan kernel $C_t$ untuk semua $t\in\F$.
\end{latihan}

\begin{latihan}\label{rank:lat-cancel}
    Misalkan $\rank(AB)=\rank B$. Buktikan bahwa, untuk setiap
    $X,Y$ berukuran sesuai,
    \[
        ABX=ABY\iff BX=BY.
    \]
    Buktikan pula versi kanan: jika $\rank(AB)=\rank A$, maka
    $XAB=YAB\iff XA=YA$.
\end{latihan}

\begin{latihan}\label{rank:lat-produk-identitas}
    Untuk $A_1,\ldots,A_k\in\Mat_n(\F)$, buktikan
    \[
        \rank(A_1A_2\cdots A_k-I)
        \le\sum_{j=1}^k\rank(A_j-I).
    \]
    Untuk kasus $A_1=\cdots=A_k=A$, temukan batas atas yang lebih
    kuat dengan memfaktorkan $A^k-I$.
\end{latihan}

\begin{latihan}\label{rank:lat-konstruksi}
    Diberikan $A\in\Mat_{m\times n}(\F)$ dengan rank $r$.
    Untuk setiap bilangan bulat $k$ dengan $r\le k\le n$, buktikan
    bahwa terdapat $B\in\Mat_n(\F)$ sedemikian sehingga
    \[
        AB=0,\qquad \rank A+\rank B=k.
    \]
\end{latihan}

\begin{latihan}\label{rank:lat-kernel-sama}
    Misalkan $A,B\in\Mat_n(\F)$. Buktikan bahwa $\ker A=\ker B$
    jika dan hanya jika terdapat matriks invertibel $P$ sehingga $A=PB$.
    Apakah kesamaan $\rank A=\rank B$ saja sudah cukup untuk
    menjamin keberadaan $P$ tersebut?
\end{latihan}

\begin{latihan}\label{rank:lat-gabungan}
    Benar atau salah: untuk setiap $A,B\in\Mat_n(\F)$,
    \[
        \rank\begin{pmatrix}A\\B\end{pmatrix}=\rank[A\ B]?
    \]
    Jika salah, berikan contoh eksplisit terkecil yang memungkinkan.
\end{latihan}

\begin{latihan}\label{rank:lat-tiga-ruang}
    Untuk matriks $A,B,C$ dengan jumlah baris sama, buktikan
    \[
        \rank[A\ B\ C]\le\rank[A\ B]+\rank[B\ C]-\rank B.
    \]
    Tentukan syarat tepat terjadinya kesamaan dalam bahasa ruang kolom.
\end{latihan}

\begin{latihan}\label{rank:lat-proyeksi}
    Untuk $0\ne x\in\C^n$, definisikan
    \[
        P=I_n-\frac{xx^*}{x^*x}.
    \]
    Tentukan rank, ruang kolom, dan kernel $P$.
\end{latihan}

\begin{latihan}\label{rank:lat-stabil}
    Buktikan bahwa $\rank A^2=\rank A$ jika dan hanya jika terdapat
    matriks invertibel $B$ dengan $A^2=BA$.
\end{latihan}

\begin{latihan}\label{rank:lat-dua-idempoten}
    Misalkan $E^2=E$ dan $F^2=F$. Buktikan bahwa $E+F$ idempoten
    jika dan hanya jika $EF=FE=0$. Jika kondisi ini dipenuhi,
    buktikan pula $\rank(E+F)=\rank E+\rank F$.
\end{latihan}

\begin{latihan}\label{rank:lat-blok-identitas}
    Untuk $A,B\in\Mat_n(\F)$, buktikan
    \[
        \rank\begin{pmatrix}A&I_n\\0&B\end{pmatrix}=n+\rank(BA).
    \]
    Gunakan hasil tersebut untuk membuktikan ketaksamaan Sylvester
    dengan bantuan \autoref{rank:segitiga-blok}.
\end{latihan}

\begin{latihan}\label{rank:lat-nilpoten}
    Diberikan $A\in\Mat_{11}(\F)$ dengan $A^3=0$ dan $\rank A=7$.
    Tentukan $\rank A^2$, $\nul A$, dan $\nul A^2$, lalu konstruksikan
    satu matriks yang memenuhi seluruh syarat tersebut.
\end{latihan}

\begin{latihan}\label{rank:lat-irisan-kernel}
    Matriks $A$ dan $B$ sama-sama mempunyai $n$ kolom, dengan
    $\nul A=a$ dan $\nul B=b$.
    Buktikan bahwa
    \[
        \max\{0,a+b-n\}\le\dim(\ker A\cap\ker B)\le\min\{a,b\}.
    \]
    Nyatakan dimensi irisan tersebut sebagai nulitas sebuah matriks blok,
    lalu tentukan syarat agar $Ax=0$ dan $Bx=0$ hanya mempunyai solusi nol.
\end{latihan}

\begin{latihan}\label{rank:lat-nulitas-produk}
    Matriks $A$ berukuran $3\times5$ dengan $\nul A=3$, sedangkan
    $B$ berukuran $5\times4$ dengan $\nul B=1$.
    Tentukan semua nilai yang mungkin bagi $\nul(AB)$ dan berikan
    konstruksi untuk setiap nilai tersebut.
\end{latihan}

\ifpembahasan
\section{Pembahasan Latihan}
\label{rank:pembahasan}

\subsection*{Latihan \ref{rank:lat-siklus}: relasi kolom sepanjang siklus}
Persamaan $C_tc=0$ memberikan
\[
    c_1-tc_n=0,\qquad -c_{j-1}+c_j=0\quad(2\le j\le n).
\]
Jadi semua $c_j$ sama, dan $(1-t)c_1=0$.
Untuk $t\ne1$, kernelnya nol dan ranknya $n$.
Untuk $t=1$, kernelnya $\operatorname{span}\{(1,\ldots,1)^T\}$
dan ranknya $n-1$. Hubungan kolom yang menyebabkan penurunan rank
adalah jumlah seluruh kolom sama dengan nol.

\subsection*{Latihan \ref{rank:lat-cancel}: pencoretan pada ruang kolom}
Dari \autoref{rank:kesamaan-produk},
$\mathcal C(B)\cap\ker A=\{0\}$.
Jika $AB(X-Y)=0$, setiap kolom $B(X-Y)$ terletak pada irisan tersebut,
sehingga $B(X-Y)=0$. Arah sebaliknya langsung.
Untuk versi kanan, transpose persamaan dan gunakan
$\rank(B^TA^T)=\rank A^T$. Pencoretan ini sah pada ruang yang relevan
walaupun $A$ atau $B$ tidak invertibel pada seluruh ruang.

\subsection*{Latihan \ref{rank:lat-produk-identitas}: ekspansi teleskopik}
Gunakan identitas
\[
    A_1\cdots A_k-I
    =\sum_{j=1}^k A_1\cdots A_{j-1}(A_j-I),
\]
dengan hasil kali kosong untuk $j=1$ diartikan $I$.
Subaditivitas rank dan batas rank hasil kali memberi
\[
    \rank(A_1\cdots A_k-I)
    \le\sum_{j=1}^k\rank\bigl(A_1\cdots A_{j-1}(A_j-I)\bigr)
    \le\sum_{j=1}^k\rank(A_j-I).
\]
Jika semua $A_j=A$, batas umum tersebut dapat diperkuat:
\[
    A^k-I=(A-I)(I+A+\cdots+A^{k-1})
    \quad\Longrightarrow\quad
    \rank(A^k-I)\le\rank(A-I).
\]
Kesamaan dapat terjadi, misalnya pada $A=0$.
Identitas khusus untuk faktor-faktor yang sama memberi informasi
lebih kuat daripada hanya menjumlahkan batas rank setiap suku.

\subsection*{Latihan \ref{rank:lat-konstruksi}: memilih citra di dalam kernel}
Karena $\dim\ker A=n-r\ge k-r$, pilih $k-r$ vektor bebas
$u_1,\ldots,u_{k-r}$ di dalam $\ker A$.
Ambil $B$ yang kolom pertamanya adalah vektor-vektor itu dan semua
kolom sisanya nol. Maka $AB=0$ dan $\rank B=k-r$.
Untuk $k=r$, cukup pilih $B=0$.
Konstruksi ini juga menjelaskan mengapa syarat $k\le n$ diperlukan.

\subsection*{Latihan \ref{rank:lat-kernel-sama}: memperluas isomorfisme}
Jika $A=PB$ dengan $P$ invertibel, jelas $Ax=0\iff Bx=0$.
Sebaliknya, asumsikan kernelnya sama dan definisikan
\[
    T:\mathcal C(B)\longrightarrow\mathcal C(A),\qquad T(Bx)=Ax.
\]
Pemetaan ini terdefinisi baik: $Bx=By$ mengakibatkan
$x-y\in\ker B=\ker A$, sehingga $Ax=Ay$.
Ia surjektif menurut definisi ruang kolom dan injektif oleh kesamaan kernel.
Pilih basis ruang kolom $B$, ambil bayangannya sebagai basis ruang
kolom $A$, lalu perluas keduanya menjadi basis $\F^n$.
Pemetaan yang memasangkan kedua basis lengkap memberikan matriks
invertibel $P$ dengan $PBx=Ax$ untuk semua $x$.

Kesamaan rank saja tidak cukup. Matriks
$A=\diag(1,0)$ dan $B=\diag(0,1)$ sama-sama mempunyai rank satu tetapi
mempunyai kernel berbeda, sehingga tidak dapat memenuhi $A=PB$
dengan $P$ invertibel.

\subsection*{Latihan \ref{rank:lat-gabungan}: ruang baris dan kolom berbeda}
Pernyataannya salah. Ambil
\[
    A=\begin{pmatrix}1&0\\0&0\end{pmatrix},\qquad
    B=\begin{pmatrix}0&1\\0&0\end{pmatrix}.
\]
Gabungan horizontal mempunyai rank satu, sedangkan gabungan vertikal
mempunyai rank dua. Ukuran $2$ adalah yang terkecil: pada ukuran $1$,
kedua gabungan sama-sama mempunyai rank nol jika kedua skalar nol,
dan rank satu jika setidaknya satu skalar tak nol.

\subsection*{Latihan \ref{rank:lat-tiga-ruang}: satu irisan yang pasti ada}
Tuliskan
$U=\mathcal C(A)+\mathcal C(B)$ dan
$V=\mathcal C(B)+\mathcal C(C)$.
Karena $\mathcal C(B)\subseteq U\cap V$, diperoleh
\begin{align*}
    \rank[A\ B\ C]
    &=\dim(U+V)=\dim U+\dim V-\dim(U\cap V)\\
    &\le\rank[A\ B]+\rank[B\ C]-\rank B.
\end{align*}
Kesamaan berlaku tepat ketika
\[
    \bigl(\mathcal C(A)+\mathcal C(B)\bigr)
    \cap\bigl(\mathcal C(B)+\mathcal C(C)\bigr)=\mathcal C(B).
\]

\subsection*{Latihan \ref{rank:lat-proyeksi}: satu arah dihilangkan}
Karena $x^*x>0$, rumus $P$ terdefinisi. Untuk setiap $y$,
$Py=y-x(x^*y)/(x^*x)$.
Jadi $Py=0$ tepat ketika $y$ kelipatan $x$.
Selanjutnya, $x^*Py=0$; sebaliknya, jika $x^*z=0$, maka $Pz=z$.
Dengan demikian
\[
    \ker P=\operatorname{span}\{x\},\qquad
    \mathcal C(P)=\{z:x^*z=0\},\qquad \rank P=n-1.
\]

\subsection*{Latihan \ref{rank:lat-stabil}: faktor invertibel dari dekomposisi}
Jika $A^2=BA$ dengan $B$ invertibel, perkalian di kiri tidak
mengubah rank, sehingga $\rank A^2=\rank A$.
Sebaliknya, \autoref{rank:stabil-awal} memberi
\[
    A=S\begin{pmatrix}D&0\\0&0\end{pmatrix}S^{-1}
\]
dengan $D$ invertibel. Pilih
$B=S\diag(D,I_{n-\rank A})S^{-1}$.
Matriks ini invertibel dan $BA=A^2$.
Untuk $A=0$, dapat langsung dipilih $B=I_n$.

\subsection*{Latihan \ref{rank:lat-dua-idempoten}: memaksa dua hasil kali nol}
Dari $(E+F)^2=E+F$ diperoleh $EF+FE=0$.
Kalikan dari kiri dan kanan dengan $E$ secara terpisah:
\[
    EF+EFE=0,\qquad EFE+FE=0.
\]
Maka $EF=FE$, sehingga $2EF=0$ dan $EF=FE=0$ di atas $\R$ atau $\C$.
Arah sebaliknya diperoleh dengan mengembangkan $(E+F)^2$.

Jika $EF=FE=0$, setiap $z\in\mathcal C(E)\cap\mathcal C(F)$ memenuhi
$Ez=z$ karena $E$ idempoten, tetapi juga $Ez=0$ karena $EF=0$.
Jadi irisannya nol. Selain itu,
$(E+F)E=E$ dan $(E+F)F=F$, sehingga
$\mathcal C(E+F)=\mathcal C(E)+\mathcal C(F)$.
Rumus dimensi memberi $\rank(E+F)=\rank E+\rank F$.

\subsection*{Latihan \ref{rank:lat-blok-identitas}: menukar blok kolom}
Pertukaran kedua blok kolom mempertahankan rank dan menghasilkan
\[
    \begin{pmatrix}I&A\\B&0\end{pmatrix}.
\]
Komplemen Schur terhadap $I$ adalah $-BA$, sehingga ranknya
$n+\rank(BA)$.
Di sisi lain, \autoref{rank:segitiga-blok} memberi batas bawah
$\rank A+\rank B$ untuk matriks semula. Maka
\[
    \rank(BA)\ge\rank A+\rank B-n.
\]
Ini adalah Sylvester untuk $BA$; tukar nama $A,B$ untuk bentuk $AB$.

\subsection*{Latihan \ref{rank:lat-nilpoten}: mengekang selisih rank}
Tuliskan $s=\rank A^2$. Barisan ranknya dimulai dengan $11,7,s,0$.
Selisih-selisihnya harus tidak meningkat:
\[
    4\ge7-s\ge s\ge0.
\]
Jadi $s\ge3$ dan $s\le7/2$; karena $s$ bilangan bulat, $s=3$.
Dengan rank--nulitas, $\nul A=11-7=4$ dan $\nul A^2=11-3=8$.
Dengan $N_3$ seperti pada contoh sebelumnya dan
$N_2=\begin{pmatrix}0&1\\0&0\end{pmatrix}$, ambil
\[
    A=\diag(N_3,N_3,N_3,N_2).
\]
Matriks ini berukuran $11$, ranknya $2+2+2+1=7$,
rank kuadratnya $1+1+1=3$, dan pangkat tiganya nol.

\subsection*{Latihan \ref{rank:lat-irisan-kernel}: dua sistem homogen sekaligus}
Ambil $U=\ker A$ dan $V=\ker B$ sebagai subruang $\F^n$.
Rumus dimensi memberi
\[
    \dim(U\cap V)=a+b-\dim(U+V)\ge a+b-n.
\]
Dimensi tidak negatif, dan $U\cap V$ termuat dalam masing-masing
ruang, sehingga juga tidak melebihi $\min\{a,b\}$.
Selanjutnya,
\[
    \ker\begin{pmatrix}A\\B\end{pmatrix}=U\cap V.
\]
Jadi sistem gabungan hanya mempunyai solusi nol tepat ketika
$\nul\binom{A}{B}=0$, atau ekuivalen $\rank\binom{A}{B}=n$.
Syarat $a+b\le n$ perlu untuk itu, tetapi belum cukup;
misalnya, dua kernel yang sama-sama berupa garis yang sama di $\F^2$
tetap beririsan tak nol.

\subsection*{Latihan \ref{rank:lat-nulitas-produk}: memperhatikan ruang perantara}
Rank kedua faktor adalah $\rank A=5-3=2$ dan $\rank B=4-1=3$.
Di ruang perantara $\F^5$, ruang $\mathcal C(B)$ dan $\ker A$
sama-sama berdimensi $3$. Dimensi irisannya sedikitnya $3+3-5=1$
dan paling banyak $3$. Karena itu,
\[
    \nul(AB)=\nul B+\dim(\mathcal C(B)\cap\ker A)\in\{2,3,4\}.
\]
Untuk merealisasikan ketiganya, tetapkan
\[
    A=\begin{pmatrix}1&0&0&0&0\\0&1&0&0&0\\0&0&0&0&0\end{pmatrix}.
\]
Dengan $e_j$ basis standar $\F^5$, pilih secara berturut-turut
\[
    B_2=[e_1\ e_2\ e_3\ 0],\quad
    B_3=[e_1\ e_3\ e_4\ 0],\quad
    B_4=[e_3\ e_4\ e_5\ 0].
\]
Masing-masing mempunyai rank $3$ dan nulitas $1$.
Hasil kali $AB_2,AB_3,AB_4$ mempunyai rank $2,1,0$ dan
nulitas $2,3,4$, sesuai konstruksi yang diminta.
\fi
